\section{Reference decoder 实现：先把一层算对}

\subsection{为什么必须先写 reference decoder}

很多推理项目失败，不是因为没有 SIMD 或 GPU，而是因为一开始没有可信 reference。后端越快，越难调；量化越低，越难判断误差来源；KV cache 越复杂，越容易把位置、head 映射和缓存槽混在一起。reference decoder 的作用是把数学合同写成最朴素、最容易审计的 C++20 代码，让后续所有优化都有比较对象。

\begin{keyidea}
reference decoder 不是低质量版本，而是正确性基准。它可以慢，但必须清楚、可复现、少隐藏状态、少复用复杂优化代码。
\end{keyidea}

本节目标不是实现完整 7B 框架，而是实现一个小型 decoder slice：token embedding、RMSNorm、Q/K/V linear、RoPE、causal attention、MLP、lm head 和 sampler 前 logits。这个 slice 可以用极小维度，例如 \code{H=16}、\code{heads=2}、\code{head\_dim=8}、\code{layers=1}，也可以用真实 shape 的小权重 fixture。重点是每一步都能输出中间张量并被 oracle 比较。

为了让每一步都能手算，本节使用固定的 TinyTrace-16：

\begin{lstlisting}[numbers=none,caption={TinyTrace-16 decoder 参数}]
vocab = 16
hidden H = 8
layers L = 2
query_heads = 2
kv_heads = 1
head_dim = 4
intermediate I = 16
prompt token ids = [3, 1, 4, 2]
\end{lstlisting}

真实 7B 模型只是把这些维度放大：\code{H} 从 8 变 4096，\code{L} 从 2 变 32，\code{I} 从 16 变 11008，\code{vocab} 从 16 变几万。reference decoder 的价值在于：TinyTrace 上每一步能手算，真实 shape 上同一套 shape 合同仍然成立。

\subsection{reference 的边界}

reference decoder 应该故意避开复杂优化：

\begin{itemize}
  \item 不做 packed weight，直接读普通 row-major 权重。
  \item 不做 SIMD intrinsic，使用普通循环。
  \item 不做线程并行，避免浮点归约顺序变化。
  \item 不复用优化后端的 int4 解包或 GPU kernel。
  \item 不把多个算子融合成一个不可观察的大函数。
\end{itemize}

它应该保留的能力是检查和观察：

\begin{itemize}
  \item 每个张量都有 shape、dtype 和 layout。
  \item 每个阶段能选择性 dump 中间结果。
  \item 每个算子能独立运行 fixture。
  \item 每个错误能指出哪一层、哪一个 head、哪一个 position。
  \item 固定 seed 的 sampler 能复现。
\end{itemize}

\subsection{最小类型：先让 shape 说话}

不要让裸指针和裸长度在 reference 中到处流动。最小实现也需要一个张量视图。

\begin{lstlisting}[numbers=none,caption={reference TensorView 的语义字段}]
TensorViewF32:
  data: span<float>
  shape: vector<int64>
  stride: vector<int64>
  name: string_view

requirements:
  shape.size == stride.size
  every index is checked in debug mode
  view does not own memory
  owner outlives view
\end{lstlisting}

reference 里可以接受 debug 检查稍慢，因为它换来定位能力。优化后端可以在 verifier 通过后走 unchecked fast path，但 reference 不应该这样做。每次 oracle 报错时，张量名、shape 和索引比一个段错误有用得多。

模型参数也要结构化，不要用字符串 map 在热路径查。

\begin{lstlisting}[numbers=none,caption={一层 decoder 的 reference 权重集合}]
DecoderLayerWeights:
  rms_attn_weight: [H]
  wq: [Q, H]
  wk: [KV, H]
  wv: [KV, H]
  wo: [H, Q]
  rms_mlp_weight: [H]
  w_gate: [I, H]
  w_up: [I, H]
  w_down: [H, I]

ModelWeights:
  token_embedding: [V, H]
  layers: array<DecoderLayerWeights>
  final_norm_weight: [H]
  lm_head: [V, H]
\end{lstlisting}

这里使用 \code{[out,in]} 形式表达线性层权重，和前面 matvec 章节一致：输出通道是行，输入通道是列。若模型文件来自别的布局，导入阶段应该转换或记录 layout，不要让 reference 每次猜。

\subsection{RMSNorm：第一个数值合同}

RMSNorm 的公式是：

\[
y_i = x_i \cdot \frac{1}{\sqrt{\frac{1}{H}\sum_j x_j^2+\epsilon}} \cdot w_i
\]

reference 实现要固定三件事：累加精度、epsilon、输出 dtype。建议 reference 用 fp32 累加，即使真实后端用 fp16/bf16。原因是 reference 要更接近数学定义，而不是模仿某个硬件的舍入。

\begin{lstlisting}[numbers=none,caption={RMSNorm reference 伪代码}]
sum_sq = 0
for i in 0..H:
  sum_sq += x[i] * x[i]

scale = rsqrt(sum_sq / H + epsilon)

for i in 0..H:
  y[i] = x[i] * scale * weight[i]
\end{lstlisting}

RMSNorm 的测试不要只测随机数。至少包含全零输入、极小值、大值、负值和 \code{H} 不能被 SIMD 宽度整除的形状。虽然 reference 不用 SIMD，但后续 optimized path 会用这些 fixture。

\subsection{Linear：所有大算子的共同骨架}

decoder 中多数重计算都是 linear。reference 版本应尽量朴素：

\begin{lstlisting}[numbers=none,caption={Linear reference 语义}]
for out in 0..OUT:
  acc = bias_or_zero(out)
  for in in 0..IN:
    acc += weight[out, in] * x[in]
  y[out] = acc
\end{lstlisting}

这段代码要配 shape 检查：

\begin{lstlisting}[numbers=none,caption={Linear shape 合同}]
x.shape == [IN]
weight.shape == [OUT, IN]
y.shape == [OUT]
if bias exists:
  bias.shape == [OUT]
\end{lstlisting}

后续 CPU 后端的 packed weight、int8、int4、SIMD micro-kernel 都要和这个 reference 比较。若 reference 自己也绕过 shape 或复用 packed descriptor，oracle 就失去意义。

\subsection{RoPE：位置编号必须作为参数传入}

RoPE 不能从数组下标偷偷推导位置。reference 函数应该显式接收 \code{model\_position}：

\begin{lstlisting}[numbers=none,caption={RoPE reference 合同}]
apply_rope(q_or_k, head_count, head_dim, model_position, rope_config):
  for head in 0..head_count:
    for pair in 0..head_dim/2:
      theta = frequency(pair, rope_config)
      angle = model_position * theta
      rotate two adjacent dimensions by angle
\end{lstlisting}

测试要覆盖下面几类错误：

\begin{itemize}
  \item position 为 0 时旋转应保持不变。
  \item 同一向量在不同 position 输出不同。
  \item prefix 长度为 \code{P} 时，新 token 的 position 是 \code{P}，不是 0。
  \item 滑动窗口或分页 cache 改变物理槽时，position 不随槽回绕。
\end{itemize}

RoPE 的 oracle 不只比较最终 logits。最好在 tiny fixture 里直接比较旋转后的 Q/K。这样一旦长上下文漂移，可以更早定位到位置合同。

\subsection{KV cache：reference 也要严格建模状态}

即使 reference 很慢，也不能每步重新计算所有历史 K/V 后假装有 KV cache。那样无法测试真实运行时状态。reference 的 KV cache 可以用普通连续数组：

\begin{lstlisting}[numbers=none,caption={KV cache reference 字段}]
KVCache:
  layers
  max_context
  kv_heads
  head_dim
  key:   [layers, max_context, kv_heads, head_dim]
  value: [layers, max_context, kv_heads, head_dim]
  filled_tokens
\end{lstlisting}

append 合同必须检查：

\begin{itemize}
  \item 写入 position 小于 \code{max\_context}。
  \item K/V 的 head 数是 \code{kv\_heads}，不是 query head 数。
  \item 写入后 \code{filled\_tokens} 和 session position 一致。
  \item debug 模式下不能覆盖已经写入但未被窗口策略淘汰的槽。
\end{itemize}

GQA 读取时要从 query head 映射到 KV head：

\begin{lstlisting}[numbers=none,caption={GQA head 映射}]
group_size = query_heads / kv_heads
kv_head = query_head / group_size
\end{lstlisting}

这个公式必须由 verifier 保证整除。若 \code{query\_heads} 不能整除 \code{kv\_heads}，模型资产应在加载阶段被拒绝，而不是让 attention 函数自己补救。

\subsection{Causal attention：先写能看懂的版本}

decode 阶段的单 token attention 最适合作为 reference 起点：

\begin{lstlisting}[numbers=none,caption={单 token decode attention reference}]
for q_head in 0..query_heads:
  kv_head = map_query_to_kv(q_head)

  for pos in 0..current_position:
    score[pos] = dot(q[q_head], key_cache[pos, kv_head]) / sqrt(head_dim)

  prob = softmax(score[0..current_position])

  out[q_head] = 0
  for pos in 0..current_position:
    out[q_head] += prob[pos] * value_cache[pos, kv_head]
\end{lstlisting}

softmax 要做数值稳定处理：

\begin{lstlisting}[numbers=none,caption={稳定 softmax}]
max_score = max(score)
sum = 0
for i:
  e[i] = exp(score[i] - max_score)
  sum += e[i]
for i:
  p[i] = e[i] / sum
\end{lstlisting}

prefill reference 可以先用同一函数逐位置执行，再逐步加批量版本。这样容易证明 causal mask：第 \code{i} 个 prompt token 只能看 \code{0..i}。后续优化可以把 prefill attention 做成矩阵形式或 fused kernel，但 oracle 仍然回到这个语义。

\subsection{MLP：SwiGLU 的中间值不能丢}

SwiGLU 的 reference 要保留 gate、up 和乘积中间值：

\begin{lstlisting}[numbers=none,caption={SwiGLU MLP reference}]
gate = linear(w_gate, x)
up = linear(w_up, x)

for i in 0..I:
  hidden[i] = silu(gate[i]) * up[i]

down = linear(w_down, hidden)
\end{lstlisting}

\code{silu(x)=x/(1+exp(-x))}。这里的中间值很重要，因为 MLP 往往是 FLOPs 大头。若最终 logits 漂移，保留 gate/up/down 的 oracle checkpoint 能迅速判断是第一段线性、激活、逐元素乘，还是 down projection 出错。

\subsection{单层 forward：把状态边写出来}

一层 decoder forward 可以写成下面的形态：

\begin{lstlisting}[numbers=none,caption={单层 decoder reference forward}]
forward_layer(layer_id, x, model_position, kv_cache):
  h = rms_norm(x, rms_attn_weight)

  q = linear(wq, h)
  k = linear(wk, h)
  v = linear(wv, h)

  apply_rope(q, query_heads, head_dim, model_position)
  apply_rope(k, kv_heads, head_dim, model_position)

  kv_cache.append(layer_id, model_position, k, v)

  attn = attention_decode(q, kv_cache, layer_id, model_position)
  x = x + linear(wo, attn)

  h = rms_norm(x, rms_mlp_weight)
  mlp = swiglu_mlp(h)
  x = x + mlp

  return x
\end{lstlisting}

注意 \code{kv\_cache.append} 是状态写入，\code{attention\_decode} 是状态读取。reference 也要把这两步分开记录。若后续 memory planner 错误复用 buffer，或者 decode step 提前读取未写完的 K/V，trace 才能定位。

\subsection{TinyTrace 的完整 shape trace}

下面只跟踪一个 token 的一层 forward。设当前 token 是 prompt 中第 2 个 token，\code{model\_position=2}，进入第 0 层前 hidden 形状是 \code{[H]=[8]}。

\begin{center}
\begin{tabular}{p{0.24\linewidth}p{0.25\linewidth}p{0.34\linewidth}}
\toprule
阶段 & 输出 shape & 必须检查的合同 \\
\midrule
embedding & \code{[8]} & token id 在 \code{[0,16)} 内 \\
RMSNorm(attn) & \code{[8]} & weight 是 \code{[8]}，epsilon 固定 \\
Q linear & \code{[8]} & \code{Wq=[8,8]} \\
K linear & \code{[4]} & \code{Wk=[4,8]}，因为 \code{kv\_heads*head\_dim=4} \\
V linear & \code{[4]} & \code{Wv=[4,8]} \\
RoPE(Q) & \code{[2,4]} & 使用 \code{position=2}，不是 cache slot \\
RoPE(K) & \code{[1,4]} & KV head 数不同于 query head 数 \\
KV append & \code{cache[0,2,0,:]} & 写第 0 层、逻辑位置 2、KV head 0 \\
attention & \code{[2,4] -> [8]} & 每个 query head 读 \code{pos 0..2} \\
O linear & \code{[8]} & \code{Wo=[8,8]} \\
RMSNorm(MLP) & \code{[8]} & residual 后再 norm \\
gate/up & \code{[16]}, \code{[16]} & \code{Wgate/Wup=[16,8]} \\
SwiGLU hidden & \code{[16]} & \code{silu(gate)*up} 逐元素 \\
down & \code{[8]} & \code{Wdown=[8,16]} \\
layer output & \code{[8]} & 第二次 residual \\
\bottomrule
\end{tabular}
\end{center}

跑完整个 prompt 后，最后一个 token 的 hidden 进入 final norm 和 lm head：

\begin{lstlisting}[numbers=none,caption={从最后 hidden 到 logits 的 shape trace}]
last_hidden: [8]
final_norm_weight: [8]
normed: [8]
lm_head weight: [16, 8]
logits: [16]
predicted_token = argmax(logits)
\end{lstlisting}

这条 trace 是 reference decoder 的最低验收线。若某个实现只告诉你“logits shape 是 \code{[16]}”，但不能说明第 0 层 position 2 的 K 写到了哪个 cache 坐标、query head 1 映射到哪个 KV head，它还不是可调试的 decoder。

\subsection{把 shape trace 绑定到仓库代码}

本仓库的 reference slice 落在 \filepath{books/cpu-volume-3/labs/linux_cpu_inference/src/reference_decoder.cpp}。它不是完整生产 runtime，但已经覆盖了这条链路的关键函数：

\begin{itemize}
  \item \code{rms\_norm}：用 fp32 累加定义 norm 合同。
  \item \code{linear\_f32}：定义 row-major \code{[out,in]} 权重语义。
  \item \code{apply\_rope}：显式接收 \code{model\_position}。
  \item \code{ReferenceKVCache::append/key/value}：按 \code{layer,position,kv\_head,dim} 寻址。
  \item \code{attention\_decode}：实现 GQA head 映射和稳定 softmax。
  \item \code{run\_reference\_decode}：从 token ids 跑到 logits 和 argmax。
  \item \code{run\_reference\_decode\_with\_trace}：输出逐 checkpoint 的 shape、stride、dtype、layout、checksum、max abs 和 values 摘要。
\end{itemize}

测试入口在 \filepath{books/cpu-volume-3/labs/linux_cpu_inference/tests/lcqi_tests.cpp}。目前它覆盖 RMSNorm、RoPE、KV attention、reference decode end-to-end、reference trace contract、packed int8 kernel 与 scalar 对齐。trace contract 会检查关键 checkpoint 的 shape、stride、dtype、layout 和 logits golden，避免 reference decoder 只验证最终 token。

当前仓库已经把 tiny reference decoder 的最终 logits 固化成 golden，而不是只检查输出范围。固定输入 \code{tokens=[1,2,3]}，\code{make\_tiny\_reference\_decoder\_model()} 的输出是：

\begin{lstlisting}[numbers=none,caption={lcqi reference decoder golden logits}]
predicted_token = 0
logits[0] =  0.352516979
logits[1] = -0.126884326
logits[2] =  0.0797837451
logits[3] = -0.29797405
logits[4] =  0.127030417
\end{lstlisting}

对应测试常量在 \code{lcqi\_tests.cpp} 中记录为 \code{LCQI\_REFERENCE\_DECODER\_LOGITS}。这仍然不是完整 oracle，但它比“predicted token 在 vocab 范围内”强得多：任何 RMSNorm、RoPE、KV append/read、attention、SwiGLU 或 lm head 的语义漂移，都会先撞到这个 logits golden。

\subsection{逐张量 checkpoint 已经能导出}

最终 logits 只能说明端到端错了，不能说明错在哪里。当前 lab 新增了 \code{lcqi\_trace} 命令，固定 prompt fixture \code{tok1 tok2 tok3}，tokenizer 输出 \code{tokens=[1,2,3]}，然后导出逐 checkpoint CSV：

\begin{lstlisting}[numbers=none,caption={reference trace 复现命令}]
books/cpu-volume-3/build/lcqi-release/labs/linux_cpu_inference/lcqi_trace \
  > books/cpu-volume-3/results/lcqi-reference-trace-2026-06-23.csv
\end{lstlisting}

CSV 字段如下：

\begin{lstlisting}[numbers=none,caption={TinyTrace checkpoint 记录}]
name,token_position,layer_id,shape,stride,dtype,layout,checksum,max_abs,values
\end{lstlisting}

前几行真实输出如下：

\begin{lstlisting}[numbers=none,caption={lcqi reference trace CSV 摘要}]
prompt,0,-1,scalar,scalar,utf8,text,0,0,tok1 tok2 tok3
tokenizer,0,-1,3,1,i32,contiguous,6,3,1;2;3
embedding,0,-1,4,1,f32,row_major,-0.25,0.4,-0.4;0.1;0.25;-0.2
attn_norm,0,0,4,1,f32,contiguous,-0.900291,1.53241,-1.53241;0.344792;1.05353;-0.766205
q_before_rope,0,0,2x2,2x1,f32,contiguous,-0.769078,0.518146,-0.518146;0.459723;-0.367778;-0.342877
k_before_rope,0,0,1x2,2x1,f32,contiguous,0.317017,0.544963,-0.227946;0.544963
kv_cache_key_slot,0,0,1x1x1x2,2x2x2x1,f32,kv_contiguous,0.317017,0.544963,-0.227946;0.544963
\end{lstlisting}

这里已经包含 shape、stride、dtype、layout 和少量 values。比如 \code{q\_before\_rope} 的 shape 是 \code{2x2}，stride 是 \code{2x1}，说明它按 \code{[query\_head,head\_dim]} 连续布局解释；\code{kv\_cache\_key\_slot} 的 shape 是 \code{1x1x1x2}，对应 \code{[layer,position,kv\_head,dim]} 的单槽切片，layout 显式标成 \code{kv\_contiguous}。

CSV 末尾会把 logits 接到 sampler 和新 token：

\begin{lstlisting}[numbers=none,caption={lcqi reference trace 的输出尾部}]
final_norm,2,-1,4,1,f32,contiguous,0.760889,1.72225,1.72225;-0.745899;0.342036;-0.557496
logits,2,-1,5,1,f32,contiguous,0.134473,0.352517,0.352517;-0.126884;0.0797838;-0.297974;0.12703
sampler_argmax,2,-1,5,1,f32,argmax,0,0,token=0
generated_token,3,-1,scalar,scalar,i32,generated,0,0,0
\end{lstlisting}

这就把 \code{prompt -> tokenizer -> embedding -> RMSNorm -> QKV -> RoPE -> attention -> KV cache -> MLP -> logits -> sampler -> token} 放进同一个可复现实验。它仍然是 tiny fixture，不是生产 tokenizer 或采样器；但它已经足够用来训练读者逐步定位 shape、layout 和数值漂移。

例如 position 2、layer 0 的关键 checkpoint：

\begin{center}
\begin{tabular}{p{0.24\linewidth}p{0.22\linewidth}p{0.34\linewidth}}
\toprule
checkpoint & shape & 定位能力 \\
\midrule
\code{emb} & \code{[8]} & token id 和 embedding 行 \\
\code{attn\_norm} & \code{[8]} & RMSNorm epsilon 和权重 \\
\code{q\_before\_rope} & \code{[2,4]} & Q projection 和 reshape \\
\code{k\_after\_rope} & \code{[1,4]} & K projection、RoPE position \\
\code{kv\_slot} & \code{[layer,pos,kv\_head]} & cache 地址和 GQA head \\
\code{attn\_prob\_h1} & \code{[3]} & causal mask 和 softmax \\
\code{mlp\_gate} & \code{[16]} & gate projection \\
\code{mlp\_hidden} & \code{[16]} & silu 和逐元素乘 \\
\code{layer\_out} & \code{[8]} & residual 和 buffer 覆盖 \\
\bottomrule
\end{tabular}
\end{center}

这些 checkpoint 不要求每次 CI 都存完整向量。常跑测试可以只比 \code{checksum/max\_abs/topk}；漂移调试时再打开完整 dump。关键是 checkpoint 名称和 shape 要稳定，否则两次报告无法对齐。

\subsection{误差边界要随 checkpoint 分层}

不同 checkpoint 对误差的容忍度不同。reference 对 reference 应该接近完全一致；优化后端若只改变循环顺序，可以用很小阈值；量化后端必须看更高层语义。

\begin{center}
\begin{tabular}{p{0.22\linewidth}p{0.24\linewidth}p{0.34\linewidth}}
\toprule
比较对象 & 推荐阈值 & 失败解释 \\
\midrule
f32 reference vs f32 scalar & \code{1e-5} abs/rel & shape、权重布局、代码错误 \\
f32 scalar vs SIMD f32 & \code{1e-4} abs/rel & 归约顺序或 SIMD lane 错 \\
int8 operator & profile dependent & scale、zero point、累加策略 \\
RoPE checkpoint & \code{1e-5} abs/rel & position、pairing、base 错 \\
attention prob & \code{1e-4} abs/rel & mask、softmax 稳定性、KV 映射 \\
logits top-k & overlap gate & 量化质量和可见输出漂移 \\
\bottomrule
\end{tabular}
\end{center}

不要用一个全局 \code{epsilon} 盖住所有层。RMSNorm 的小误差可能还能接受；attention score 的小误差经过 softmax 后可能改变概率；logits 的小差异若正好发生在 top-1/top-2 附近，会直接改变采样结果。

\subsection{端到端 reference decode loop}

端到端 reference 可以先只输出 logits，不做复杂采样：

\begin{lstlisting}[numbers=none,caption={reference decode loop}]
tokens = tokenizer_fixture.encode(prompt)

for pos in 0..tokens.size:
  x = embedding[tokens[pos]]
  for layer in 0..layers:
    x = forward_layer(layer, x, pos, kv_cache)

last_hidden = x
normed = rms_norm(last_hidden, final_norm_weight)
logits = linear(lm_head, normed)
\end{lstlisting}

prefill 的最后一个 token 产生首个 logits。之后 decode 追加 sampler 选出的 token：

\begin{lstlisting}[numbers=none,caption={decode 生成循环 reference}]
for step in 0..max_new_tokens:
  next_id = sampler(logits, seed, params)
  tokens.push_back(next_id)

  pos = tokens.size - 1
  x = embedding[next_id]
  for layer in 0..layers:
    x = forward_layer(layer, x, pos, kv_cache)

  logits = linear(lm_head, rms_norm(x, final_norm_weight))
\end{lstlisting}

这里故意没有 batch、线程和 fusion。它的价值是每个 token、每层、每个中间值都能被检查。

\subsection{Golden fixture 怎么设计}

reference 必须配 golden fixture。一个高质量 fixture 至少包含：

\begin{lstlisting}[numbers=none,caption={tiny decoder fixture 内容}]
model_dims:
  vocab
  hidden
  layers
  query_heads
  kv_heads
  head_dim
  intermediate
  max_context

weights:
  token_embedding
  layer weights
  final_norm
  lm_head

input:
  prompt_token_ids
  sampler_params
  seed

expected:
  checkpoints
  logits
  generated_token_ids
\end{lstlisting}

权重不要全用随机数。建议混合使用小整数、正负数、接近零的值和非对称值。全随机 fixture 容易覆盖到计算路径，却不容易人工定位错误。比如 RoPE 可以用只有两个维度非零的向量，attention 可以用能手算 softmax 排序的 Q/K/V。

\subsection{checkpoint 策略}

不要只保存最终 logits。推荐 checkpoint 分层：

\begin{center}
\begin{tabular}{p{0.20\linewidth}p{0.36\linewidth}p{0.28\linewidth}}
\toprule
checkpoint & 目的 & 常见错误 \\
\midrule
embedding & tokenizer 和查表 & token id 错、词表错 \\
norm output & RMSNorm & epsilon、累加精度 \\
q/k/v & linear 与 head reshape & 权重转置、shape 错 \\
rope q/k & 位置编码 & position、维度配对 \\
attention prob & mask 与 softmax & 未来泄漏、溢出 \\
attention out & KV cache 读取 & head 映射、cache slot \\
mlp hidden & SwiGLU & 激活、逐元素乘 \\
layer output & residual & 加法位置、覆盖 buffer \\
logits top-k & 可见输出 & lm head、采样差异 \\
\bottomrule
\end{tabular}
\end{center}

checkpoint 太多会增加文件体积，可以通过 \code{trace\_level} 控制。日常测试比较关键 checkpoint；调试漂移时打开 debug 全量输出。

\subsection{reference 和优化后端怎样比较}

比较不能只看完全相等。不同 dtype 和不同归约顺序会造成微小差异。oracle 应按层级使用不同阈值：

\begin{lstlisting}[numbers=none,caption={oracle tolerance profile}]
f32_scalar:
  absolute: 1e-5
  relative: 1e-5

f16_or_bf16:
  absolute: 2e-2
  relative: 2e-2

int8_weight:
  per_operator_absolute: profile dependent
  logits_topk_overlap: required

int4_weight:
  compare logits, top-k and generation under fixed prompt set
\end{lstlisting}

阈值必须随 profile 记录到报告里。不要让测试代码里一个魔法数决定所有后端。特别是量化后端，operator 误差、logits top-k 和生成序列要一起看；单个最终文本“看起来不错”不能证明模型没有漂移。

\subsection{本节验收线}

读完本节，应该能动手写一个小型 reference decoder：

\begin{itemize}
  \item 用 \code{TensorView} 和结构化权重表达 shape，而不是裸指针乱传。
  \item 分别实现 RMSNorm、Linear、RoPE、KV append/read、attention、SwiGLU 和 lm head。
  \item 显式传入 \code{model\_position}，不把物理 cache slot 当 RoPE 位置。
  \item 用 tiny fixture 保存 checkpoint 和 logits golden。
  \item 用 oracle profile 比较 reference 与优化后端。
  \item 把 reference 当正确性基准，而不是被优化路径污染的慢版本。
\end{itemize}
